\begin{lemma}
Let $G$ be a $\Delta$-regular graph, let $r\in V(G)$, and let
$X\subseteq N_G(r)$.  In proving the sampling bound for $X$, one may replace
$G$ by a $\Delta$-regular graph $G'$ with an embedded copy of
$G[X\cup\{r\}]$ such that the embedded copy has the same adjacencies as
$G[X\cup\{r\}]$, no two distinct vertices of $X$ have a common neighbor in
$G'$ outside the closed neighborhood of $r$, and the activation--priority
law on $G'$ gives
\[
\Pr_G[I\cap X\ne\emptyset]\le
\Pr_{G'}[I'\cap X\ne\emptyset].
\]
In particular the local graph used to compute
$b_{G[X\cup\{r\}],\Delta}(r)$ is preserved by the reduction.
\end{lemma}
\begin{proof}
Write $L=X\cup\{r\}$ and first reserve an embedded copy of the induced graph
$G[L]$.  This copy is never changed.  Thus the root is adjacent in the copy
precisely to the vertices of $X$, the adjacencies and non-adjacencies inside
$X$ are exactly those of $G$, and the induced rooted graph which enters
$b_{G[X\cup\{r\}],\Delta}(r)$ is preserved.

We next build a finite graph $H$ around this fixed copy.  For each incidence
$(x,z)$ with $x\in X$, $z\in V(G)\setminus L$, and $xz\in E(G)$, add a fresh
vertex $z_x$ adjacent only to the embedded copy of $x$.  Thus an outside
vertex of $G$ adjacent to several vertices of $X$ is replaced by separate
one-neighbor blockers, one for each of its incidences with $X$.  Since $G$ is
$\Delta$-regular, the embedded vertex $x$ has
\[
\Delta-1-d_{G[X]}(x)
\]
neighbors outside $L$ in $G$, and after adding these incidence blockers it
has degree
$1+d_{G[X]}(x)+(\Delta-1-d_{G[X]}(x))=\Delta$ in $H$.  The vertex $r$ has
degree $|X|$ in the fixed copy.  If $|X|<\Delta$, add
$\Delta-|X|$ further fresh vertices adjacent only to $r$; if $|X|=\Delta$,
add none.  These extra vertices are deliberately placed outside the embedded
local graph, so the induced graph on $L$ remains $G[L]$, but $r$ now also has
degree $\Delta$.

Every vertex of $H$ has degree at most $\Delta$: the embedded vertices in
$X\cup\{r\}$ already have degree $\Delta$, and every fresh outside vertex has
degree $1$.  Complete the remaining degree deficits by a finite simple
regular filler which is disjoint from the embedded local graph except at the
vertices whose deficits are being filled.  Concretely, list one formal stub
for each missing incident edge of a vertex of $H$.  If the number of stubs is
odd, add one isolated fresh vertex; this can occur only when $\Delta$ is odd,
and the isolated vertex contributes exactly $\Delta$ additional stubs, making
the total number even.  Pair the stubs arbitrarily.  For each paired pair of
stubs, add a fresh copy of the complete graph
$K_{\Delta+1}$, delete one edge $ab$ of that copy, and join the two stubs to
$a$ and $b$, respectively.  The two endpoints $a,b$ had degree $\Delta-1$
after deleting $ab$ and regain degree $\Delta$ after the two new joining
edges; all other vertices of the copy keep degree $\Delta$; and each old
stub is filled by one new edge.  Even if the two stubs in a pair belong to
the same old vertex, the two joining edges go to the distinct fresh vertices
$a$ and $b$ of a previously unused block.  Hence no multiple edge is created
and no new edge is inserted inside the embedded copy of $L$.  The resulting
finite simple graph is $\Delta$-regular; call it $G'$.

By construction, every vertex of $G'$ outside the closed neighborhood of the
embedded root $r$ is adjacent to at most one embedded vertex of $X$.  Indeed,
the only fresh vertices adjacent to $X$ are the incidence blockers $z_x$, and
each such vertex is adjacent to exactly the one vertex $x$ of $X$; the filler
blocks attach to those blockers and to the extra neighbors of $r$, but never
to vertices of $X$ or to two vertices of $X$.  Therefore no two distinct
vertices of $X$ have a common neighbor in $G'$ outside $N_{G'}[r]$.

It remains to compare the sampling probabilities.  By
\noderef{RandomIndependentSetSampling}, the random set is generated from
independent activation variables with activation probability
$p=\gamma/\Delta$ and independent uniform priorities, and a vertex survives
exactly when it is activated and has larger priority than every activated
neighbor.  Couple the activation and priority variables on the embedded copy
of $L$ identically in $G$ and $G'$.  The extra neighbors of $r$ and all filler
vertices not adjacent to $X$ may be sampled arbitrarily according to the
product law, because they are not neighbors of vertices in $X$ and hence
cannot affect the event $I\cap X\ne\emptyset$.

Fix the variables on $L$ and expose all outside blockers except one original
outside vertex $z$ of $G$ which is adjacent to a nonempty set
$B\subseteq X$.  Let $C\subseteq B$ be the set of vertices of $B$ which,
after the already exposed variables are taken into account, would survive
provided they are not blocked by $z$ or by its replacements.  If some vertex
of $X\setminus B$ already survives, then both graphs have the event
$I\cap X\ne\emptyset$ regardless of $z$.  If $C=\emptyset$, then neither the
single blocker nor its replacements can create a survivor in $B$.  Thus the
only case to compare is $C\ne\emptyset$ and no vertex outside $B$ has already
survived.

In the original graph, $z$ kills all vertices of $C$ exactly when $z$ is
activated and has priority larger than every priority $\pi(x)$ with
$x\in C$.  Conditional on the priorities on $C$, this has probability
$p(1-\max_{x\in C}\pi(x))$.  In the split graph, the independent replacement
blockers $z_x$ kill all vertices of $C$ with probability
\[
\prod_{x\in C} p(1-\pi(x)).
\]
Since $0\le p\le 1$ and $0\le\pi(x)\le 1$, this product is at most
$p(1-\max_{x\in C}\pi(x))$: choose $x_0$ with maximal priority, keep the
factor $p(1-\pi(x_0))$, and bound every other factor by $1$.  Hence replacing
one shared outside blocker by independent incidence blockers cannot decrease,
after conditioning on all other variables, the probability that at least one
vertex of $X$ survives.

Apply this one-blocker comparison successively to every outside vertex of
$G$ adjacent to at least one vertex of $X$ but lying outside $L$.  The
successive replacements are finite, and the conditional inequality at each
step is preserved after averaging over the variables already exposed.  At
the end of the replacement process we obtain exactly the blocker interface
used in the construction of $G'$.  Therefore the activation--priority law on
$G'$ satisfies
\[
\Pr_G[I\cap X\ne\emptyset]\le
\Pr_{G'}[I'\cap X\ne\emptyset].
\]
Together with the preserved induced copy of $G[X\cup\{r\}]$, this proves the
claimed reduction and preserves the local parameter
$b_{G[X\cup\{r\}],\Delta}(r)$.
\end{proof}
